\documentclass[11pt,aspectratio=169]{beamer}

% ---------------------------------------------------------------------------
% WirSchiffenDas - Engine Quality Analysis
% Beamer-Praesentation: Technische Anforderungen und ADRs im Code
% ---------------------------------------------------------------------------

\usepackage[T1]{fontenc}
\usepackage[utf8]{inputenc}
\usepackage[ngerman]{babel}
\usepackage{lmodern}
\usepackage{microtype}
\usepackage{booktabs}
\usepackage{tabularx}
\usepackage{array}
\usepackage{makecell}
\usepackage{xcolor}
\usepackage{tikz}
\usetikzlibrary{arrows.meta,positioning,shapes.geometric,fit,backgrounds}
\usepackage{listings}
\usepackage{seqsplit}
\usepackage{hyperref}

% --- Tabellenspaltentypen (wie in architecture.tex) ------------------------
\newcolumntype{Y}{>{\raggedright\arraybackslash}X}
\newcolumntype{L}[1]{>{\raggedright\arraybackslash}p{#1}}

% --- Farben (identisch zur arc42-Dokumentation) ----------------------------
\definecolor{hbrsblue}{HTML}{00A6C8}
\definecolor{darkblue}{HTML}{0A556B}
\definecolor{lightblue}{HTML}{E8F7FA}
\definecolor{lightgray}{HTML}{F3F5F6}
\definecolor{warnorange}{HTML}{B26A00}
\definecolor{okgreen}{HTML}{2E7D32}
\definecolor{codegray}{HTML}{6B7280}
\definecolor{adryellow}{HTML}{FCBA03}

\hypersetup{
  colorlinks=true,
  linkcolor=darkblue,
  urlcolor=darkblue,
  pdftitle={WirSchiffenDas - TA und ADRs im Code}
}

% --- Theme -----------------------------------------------------------------
\usetheme{default}
\usecolortheme{default}
\setbeamertemplate{navigation symbols}{}

\setbeamercolor{structure}{fg=darkblue}
\setbeamercolor{frametitle}{fg=white,bg=darkblue}
\setbeamercolor{title}{fg=darkblue}
\setbeamercolor{subtitle}{fg=darkblue}
\setbeamercolor{author}{fg=black}
\setbeamercolor{date}{fg=black}
\setbeamercolor{normal text}{fg=black}
\setbeamercolor{itemize item}{fg=hbrsblue}
\setbeamercolor{itemize subitem}{fg=hbrsblue}
\setbeamercolor{enumerate item}{fg=hbrsblue}
\setbeamercolor{block title}{fg=white,bg=darkblue}
\setbeamercolor{block body}{bg=lightblue}
\setbeamercolor{alerted text}{fg=warnorange}
\setbeamercolor{author in head/foot}{fg=codegray,bg=lightgray}
\setbeamercolor{title in head/foot}{fg=codegray,bg=lightgray}
\setbeamercolor{date in head/foot}{fg=codegray,bg=lightgray}

\setbeamerfont{frametitle}{size=\large,series=\bfseries}
\setbeamerfont{title}{size=\Huge,series=\bfseries}
\setbeamerfont{subtitle}{size=\large}
\setbeamerfont{author in head/foot}{size=\tiny}
\setbeamerfont{title in head/foot}{size=\tiny}
\setbeamerfont{date in head/foot}{size=\tiny}

\setbeamertemplate{frametitle}{%
  \nointerlineskip%
  \begin{beamercolorbox}[wd=\paperwidth,ht=2.9ex,dp=1.3ex,leftskip=0.35cm,rightskip=0.35cm]{frametitle}%
    \usebeamerfont{frametitle}\insertframetitle%
  \end{beamercolorbox}%
}

\setbeamertemplate{footline}{%
  \leavevmode%
  \hbox{%
    \begin{beamercolorbox}[wd=.34\paperwidth,ht=2.4ex,dp=1ex,leftskip=.35cm]{author in head/foot}%
      \usebeamerfont{author in head/foot}\insertshortauthor%
    \end{beamercolorbox}%
    \begin{beamercolorbox}[wd=.42\paperwidth,ht=2.4ex,dp=1ex,center]{title in head/foot}%
      \usebeamerfont{title in head/foot}\insertshorttitle%
    \end{beamercolorbox}%
    \begin{beamercolorbox}[wd=.24\paperwidth,ht=2.4ex,dp=1ex,rightskip=.35cm]{date in head/foot}%
      \usebeamerfont{date in head/foot}\hfill\insertframenumber\,/\,\inserttotalframenumber%
    \end{beamercolorbox}%
  }%
  \vskip0pt%
}

\setbeamertemplate{itemize items}[triangle]
\setbeamertemplate{itemize subitem}[circle]

% Kompaktere Listen und Bloecke, damit Code-Snippets Platz haben
\setbeamertemplate{itemize/enumerate body begin}{\small}
\setbeamertemplate{itemize/enumerate subbody begin}{\footnotesize}
\setlength{\itemsep}{1pt}
\setlength{\parskip}{1pt}
\setbeamersize{text margin left=0.6cm,text margin right=0.6cm}

% --- Code-Highlighting -----------------------------------------------------
\lstdefinestyle{base}{
  basicstyle=\ttfamily\tiny,
  keywordstyle=\color{darkblue}\bfseries,
  commentstyle=\color{codegray}\itshape,
  stringstyle=\color{okgreen},
  numberstyle=\tiny\color{codegray},
  numbers=left,
  numbersep=5pt,
  showstringspaces=false,
  breaklines=true,
  breakatwhitespace=false,
  frame=single,
  framesep=3pt,
  rulecolor=\color{lightgray},
  backgroundcolor=\color{lightgray!45},
  tabsize=2,
  columns=fullflexible,
  keepspaces=true,
  aboveskip=3pt,
  belowskip=1pt,
}

\lstdefinestyle{java}{style=base,language=Java}
\lstdefinestyle{yaml}{style=base,language=,morekeywords={services,image,ports,volumes,environment,depends_on,resilience4j,circuitbreaker,instances}}
\lstdefinestyle{docker}{style=base,language=,morekeywords={FROM,WORKDIR,COPY,RUN,EXPOSE,ENTRYPOINT,AS}}
\lstdefinestyle{bash}{style=base,language=,morekeywords={docker,compose,up,build,stop,start,curl,POST,GET}}

% --- Hilfsmakros -----------------------------------------------------------
\newcommand{\hbrslogo}{%
\begin{tikzpicture}[baseline=-0.4ex,scale=0.6]
  \draw[line width=2.5pt,hbrsblue] (0,0) circle (0.34);
  \fill[hbrsblue] (1.0,0) circle (0.34);
  \node[anchor=west,text=black,font=\sffamily\bfseries\large] at (1.55,0.15) {Hochschule};
  \node[anchor=west,text=black,font=\sffamily\bfseries\large] at (1.55,-0.5) {Bonn-Rhein-Sieg};
\end{tikzpicture}}

\newcommand{\code}[1]{\texttt{\small\detokenize{#1}}}
\newcommand{\apiendpoint}[1]{\path{#1}}

% Badge fuer Anforderungs-/ADR-Kennung oben rechts im Frame
\newcommand{\reqbadge}[2]{%
  \begin{tikzpicture}[remember picture,overlay]
    \node[anchor=north east,xshift=-0.35cm,yshift=-0.15cm,
          fill=#1,rounded corners=2pt,inner sep=3pt,
          font=\sffamily\bfseries\scriptsize,text=white] at (current page.north east) {#2};
  \end{tikzpicture}}

\newcommand{\tabadge}[1]{\reqbadge{hbrsblue}{#1}}
\newcommand{\adrbadge}[1]{\reqbadge{adryellow}{#1}}
\newcommand{\frbadge}[1]{\reqbadge{okgreen}{#1}}
\newcommand{\qrbadge}[1]{\reqbadge{warnorange}{#1}}

% Datei-Hinweis unter einem Snippet (umbrechbar)
\newcommand{\srcfile}[1]{{\scriptsize\color{codegray}\ttfamily\seqsplit{#1}}}

% --- Metadaten -------------------------------------------------------------
\title[TA \& ADRs im Code]{Technische Anforderungen und ADRs im Code}
\subtitle{WirSchiffenDas -- Engine Quality Analysis}
\author[SEKA, SS 2026]{Theodoros Bampis \& Yunus Emre Sirin}
\date{Sommersemester 2026}

\begin{document}

% --- Titelfolie ------------------------------------------------------------
\begin{frame}[plain]
  \begin{center}
    \hbrslogo

    \vspace{1.0cm}
    {\usebeamerfont{title}\usebeamercolor[fg]{title}\inserttitle\par}
    \vspace{0.3cm}
    {\usebeamerfont{subtitle}\usebeamercolor[fg]{subtitle}\insertsubtitle\par}

    \vspace{0.9cm}
    {\large \insertauthor\par}
    \vspace{0.15cm}
    {\large \insertdate\par}

    \vspace{0.8cm}
    \begin{tikzpicture}
      \node[fill=hbrsblue!12,rounded corners=3pt,inner sep=6pt,
            font=\sffamily\bfseries\small,text=darkblue]
        {Service-basierte Komponenten-Architekturen (SEKA)};
    \end{tikzpicture}
  \end{center}
\end{frame}

% --- Inhalt -----------------------------------------------------------------
% ===========================================================================
% 01 - Einordnung: Agenda, Ziel und Traceability-Ansatz
% ===========================================================================

\begin{frame}{Agenda}
  \begin{columns}[T,onlytextwidth]
    \begin{column}{0.48\textwidth}
      \textbf{Teil 1 -- Einordnung}
      \begin{itemize}
        \item Ziel der Präsentation
        \item Traceability-Ansatz
      \end{itemize}
      \vspace{0.4cm}
      \textbf{Teil 2 -- Architektur}
      \begin{itemize}
        \item Service-Landschaft
        \item Choreographie
      \end{itemize}
    \end{column}
    \begin{column}{0.48\textwidth}
      \textbf{Teil 3 -- Technische Anforderungen}
      \begin{itemize}
        \item TA-01 bis TA-05 im Code
      \end{itemize}
      \vspace{0.4cm}
      \textbf{Teil 4 -- Architekturentscheidungen}
      \begin{itemize}
        \item ADR-01 bis ADR-06 im Code
      \end{itemize}
      \vspace{0.4cm}
      \textbf{Teil 5 -- Demo \& Fazit}
    \end{column}
  \end{columns}

  \vfill
  \begin{block}{Leitfrage}
    \small Welche Anforderung ist \textbf{wo} im Code nachweisbar -- und
    \textbf{warum} wurde sie so umgesetzt?
  \end{block}
\end{frame}

\begin{frame}{Ziel: Nachweis statt Behauptung}
  \begin{columns}[T,onlytextwidth]
    \begin{column}{0.54\textwidth}
      \textbf{Problem}
      \begin{itemize}
        \item Architekturdokumente beschreiben Anforderungen.
        \item Ohne Code-Bezug bleibt offen, ob sie umgesetzt sind.
      \end{itemize}

      \vspace{0.3cm}
      \textbf{Lösung: Traceability}
      \begin{itemize}
        \item Jede Anforderung $\rightarrow$ konkrete Datei + Symbol.
        \item Jedes Snippet $\rightarrow$ kurze Erläuterung.
        \item Nachvollziehbar in der mündlichen Prüfung.
      \end{itemize}
    \end{column}
    \begin{column}{0.42\textwidth}
      \begin{block}{Kennungen}
        \small
        \begin{tabular}{@{}ll@{}}
          \tabadge{TA-xx} & Technische Anforderung \\
          \adrbadge{ADR-xx} & Entwurfsentscheidung \\
          \frbadge{FR-xx} & Funktionale Anforderung \\
          \qrbadge{QR-xx} & Qualitätsanforderung \\
        \end{tabular}
      \end{block}
      \vspace{0.2cm}
      \begin{block}{Quelle}
        \small \code{docs/code-traceability.md}
      \end{block}
    \end{column}
  \end{columns}
\end{frame}

\begin{frame}{Traceability-Matrix (Auszug)}
  \small
  \begin{tabularx}{\textwidth}{L{1.6cm}L{4.6cm}Y}
    \toprule
    \textbf{Kennung} & \textbf{Anforderung} & \textbf{Zentraler Code-Nachweis} \\
    \midrule
    TA-01 & Spring Boot & \code{@SpringBootApplication} in 6 Service-Klassen \\
    TA-02 & Docker & 6 $\times$ \code{Dockerfile} (Multi-Stage) \\
    TA-03 & Docker Compose & \code{docker-compose.yml} (6 Services, Volumes) \\
    TA-04 & Resilience4j & \code{@CircuitBreaker} + \code{fallback} + \code{application.yml} \\
    TA-05 & Kafka (COULD) & bewusst nicht umgesetzt (ADR-02) \\
    \midrule
    ADR-01 & Fachlicher Schnitt & Modulstruktur \code{api/application/domain/infrastructure} \\
    ADR-03 & Choreographie & \code{AnalysisWorker} $\rightarrow$ \code{NextServiceClient} \\
    ADR-06 & Keine Shared Persistence & \code{ConfigurationSnapshot} + stateless Services \\
    \bottomrule
  \end{tabularx}

  \vfill
  \begin{center}
    \small\color{codegray}
    Vollständige Matrix: \code{docs/code-traceability.md}
  \end{center}
\end{frame}

% ===========================================================================
% 02 - Architektur: Service-Landschaft und Choreographie
% ===========================================================================

\begin{frame}{Service-Landschaft}
  \begin{columns}[T,onlytextwidth]
    \begin{column}{0.5\textwidth}
      \begin{itemize}
        \item \textbf{6 Microservices}, je ein Spring-Boot-Prozess.
        \item \textbf{2 zustandsbehaftet}: Configuration, Analysis Management.
        \item \textbf{4 stateless}: die Analyse-Services.
        \item Kommunikation ausschließlich über REST.
      \end{itemize}

      \vspace{0.3cm}
      \begin{block}{Fachlicher Schnitt}
        \small Jeder Service kapselt eine Business Capability --
        kein technischer Layer-Schnitt (ADR-01).
      \end{block}
    \end{column}
    \begin{column}{0.46\textwidth}
      \begin{center}
      \begin{tikzpicture}[
        node distance=4mm and 8mm,
        every node/.style={font=\scriptsize},
        svc/.style={draw=darkblue,rounded corners,fill=lightblue,minimum width=34mm,minimum height=6mm,align=center},
        data/.style={draw,ellipse,fill=lightgray,minimum width=26mm,minimum height=5mm,align=center},
        ana/.style={draw=darkblue,rounded corners,fill=white,minimum width=34mm,minimum height=6mm,align=center}
      ]
      \node[svc] (cfg) {configuration-service};
      \node[svc,below=of cfg] (am) {analysis-management-service};
      \node[data,below=of cfg] (cfgdb) {Configuration DB};
      \node[data,below=of am] (amdb) {Analysis DB};
      \node[ana,below=of amdb] (fluid) {fluid-analysis};
      \node[ana,below=of fluid] (thermal) {thermal-analysis};
      \node[ana,below=of thermal] (elec) {electrical-analysis};
      \node[ana,below=of elec] (ems) {engine-management};
      \draw[-{Latex}] (cfg) -- (cfgdb);
      \draw[-{Latex}] (am) -- (amdb);
      \draw[-{Latex},thick] (am) -- (fluid);
      \draw[-{Latex},thick] (fluid) -- (thermal);
      \draw[-{Latex},thick] (thermal) -- (elec);
      \draw[-{Latex},thick] (elec) -- (ems);
      \end{tikzpicture}
      \end{center}
    \end{column}
  \end{columns}
\end{frame}

\begin{frame}{Choreographie statt Orchestration}
  \begin{columns}[T,onlytextwidth]
    \begin{column}{0.5\textwidth}
      \textbf{Ablauf}
      \begin{enumerate}
        \item Analysis Management startet \textbf{nur} den Anchor \code{FLUID}.
        \item Jeder erfolgreiche Service ruft selbst den nächsten auf.
        \item Ergebnisse werden im \code{AnalysisCommand} mitgereicht.
        \item Engine Management berücksichtigt Vorergebnisse.
      \end{enumerate}

      \vspace{0.2cm}
      \begin{block}{Kernaussage}
        \small Kein zentraler Orchestrator -- die Ablaufsteuerung liegt
        verteilt in den Services (ADR-03).
      \end{block}
    \end{column}
    \begin{column}{0.46\textwidth}
      \begin{center}
      \begin{tikzpicture}[
        node distance=5mm,
        every node/.style={font=\scriptsize},
        stepbox/.style={draw=darkblue,rounded corners,fill=lightblue,minimum width=20mm,minimum height=8mm,align=center}
      ]
      \node[stepbox] (am) {Analysis\\Management};
      \node[stepbox,right=of am] (f) {Fluid};
      \node[stepbox,right=of f] (t) {Thermal};
      \node[stepbox,below=of t] (e) {Electrical};
      \node[stepbox,left=of e] (m) {Engine\\Management};
      \draw[-{Latex},thick] (am) -- node[above]{start} (f);
      \draw[-{Latex},thick] (f) -- (t);
      \draw[-{Latex},thick] (t) -- (e);
      \draw[-{Latex},thick] (e) -- (m);
      \end{tikzpicture}
      \end{center}

      \vspace{0.2cm}
      \begin{center}
      \begin{tikzpicture}[
        every node/.style={font=\scriptsize},
        cb/.style={draw=warnorange,rounded corners,fill=warnorange!10,minimum width=26mm,minimum height=8mm,align=center}
      ]
      \node[cb] {Circuit Breaker\\Thermal nicht erreichbar};
      \end{tikzpicture}
      \end{center}
    \end{column}
  \end{columns}
\end{frame}

% ===========================================================================
% 03 - Technische Anforderungen (TA) im Code
% ===========================================================================

% --- TA-01 -----------------------------------------------------------------
\begin{frame}[fragile]{TA-01 -- Spring Boot}
  \tabadge{TA-01}
  \begin{columns}[T,onlytextwidth]
    \begin{column}{0.52\textwidth}
      \begin{lstlisting}[style=java]
@EnableAsync
@SpringBootApplication
public class FluidAnalysisServiceApplication {
    public static void main(String[] args) {
        SpringApplication.run(
            FluidAnalysisServiceApplication.class, args);
    }
}
      \end{lstlisting}
      \srcfile{services/fluid-analysis-service/.../FluidAnalysisServiceApplication.java}
    \end{column}
    \begin{column}{0.44\textwidth}
      \textbf{Erläuterung}
      \begin{itemize}
        \item \code{@SpringBootApplication} aktiviert Auto-Configuration, Component Scan und die eingebettete Servlet-Engine.
        \item \code{@EnableAsync} ist die Grundlage für die asynchrone Ausführung (QR-08).
        \item Identisches Muster in allen sechs Services.
      \end{itemize}
      \vspace{0.2cm}
      \begin{block}{Abhängigkeit}
        \small \code{pom.xml}: \code{spring-boot-starter-web}
      \end{block}
    \end{column}
  \end{columns}
\end{frame}

% --- TA-02 -----------------------------------------------------------------
\begin{frame}[fragile]{TA-02 -- Docker (Container pro Service)}
  \tabadge{TA-02}
  \begin{columns}[T,onlytextwidth]
    \begin{column}{0.54\textwidth}
      \begin{lstlisting}[style=docker]
FROM maven:3.9.16-eclipse-temurin-21 AS build
WORKDIR /workspace
COPY pom.xml ./
COPY services ./services
RUN mvn -q -pl services/fluid-analysis-service \
    -am clean package -DskipTests

FROM eclipse-temurin:21-jre
WORKDIR /app
COPY --from=build /workspace/services/fluid-analysis-service/target/fluid-analysis-service-0.1.0-SNAPSHOT.jar app.jar
EXPOSE 8083
ENTRYPOINT ["java", "-jar", "/app/app.jar"]
      \end{lstlisting}
      \srcfile{services/fluid-analysis-service/Dockerfile}
    \end{column}
    \begin{column}{0.42\textwidth}
      \textbf{Erläuterung}
      \begin{itemize}
        \item \textbf{Multi-Stage-Build}: Build-Umgebung (Maven + JDK) getrennt vom schlanken Runtime-Image (nur JRE).
        \item Jeder Service erhält ein eigenes Image und einen eigenen Port.
        \item Grundlage für Independent Deployability (QR-03, ADR-05).
      \end{itemize}
      \vspace{0.2cm}
      \begin{block}{Nachweis}
        \small 6 $\times$ \code{Dockerfile} -- eines pro Service.
      \end{block}
    \end{column}
  \end{columns}
\end{frame}

% --- TA-03 -----------------------------------------------------------------
\begin{frame}[fragile]{TA-03 -- Docker Compose}
  \tabadge{TA-03}
  \begin{columns}[T,onlytextwidth]
    \begin{column}{0.54\textwidth}
      \begin{lstlisting}[style=yaml]
services:
  configuration-service:
    image: ysirin2s/seka-wirschaffendas:configuration-service-v${VERSION}
    ports:
      - "8081:8081"
    volumes:
      - configuration-data:/app/data
  analysis-management-service:
    image: ysirin2s/seka-wirschaffendas:analysis-management-service-v${VERSION}
    ports:
      - "8082:8082"
    environment:
      CONFIGURATION_SERVICE_URL: http://configuration-service:8081
      FLUID_ANALYSIS_SERVICE_URL: http://fluid-analysis-service:8083
    volumes:
      - analysis-data:/app/data
      \end{lstlisting}
      \srcfile{docker-compose.yml}
    \end{column}
    \begin{column}{0.42\textwidth}
      \textbf{Erläuterung}
      \begin{itemize}
        \item Compose definiert alle sechs Services, Ports, Umgebungsvariablen und getrennte Volumes.
        \item Services finden sich über \textbf{Service-Namen} im Netzwerk \code{wirschaffendas-network} -- keine fest verdrahteten IPs.
        \item \code{depends_on} steuert die Startreihenfolge.
      \end{itemize}
      \begin{block}{Start}
        \small \code{docker compose up --build}
      \end{block}
    \end{column}
  \end{columns}
\end{frame}

% --- TA-04 -----------------------------------------------------------------
\begin{frame}[fragile]{TA-04 -- Resilience4j: Circuit Breaker}
  \tabadge{TA-04}
  \begin{columns}[T,onlytextwidth]
    \begin{column}{0.56\textwidth}
      \begin{lstlisting}[style=java]
@CircuitBreaker(name = "nextService", fallbackMethod = "fallback")
public void startNext(AnalysisCommand command, String currentResult) {
    var results = new ArrayList<>(command.previousResults());
    results.add(Map.of("algorithm", "FLUID", "result", currentResult));
    AnalysisCommand nextCommand = new AnalysisCommand(
        command.analysisId(), command.configuration(), results);
    client.post().uri("/internal/analyses")
        .body(nextCommand).retrieve().toBodilessEntity();
}
private void fallback(AnalysisCommand command,
                      String currentResult, Throwable throwable) {
    managementClient.reportStatus(command.analysisId(), "THERMAL",
        "FAILED", "thermal-analysis-service unavailable");
}
      \end{lstlisting}
      \srcfile{services/fluid-analysis-service/.../infrastructure/NextServiceClient.java}
    \end{column}
    \begin{column}{0.4\textwidth}
      \textbf{Erläuterung}
      \begin{itemize}
        \item Der Breaker fängt Verbindungsfehler ab.
        \item Statt Fehlerkaskade läuft der \textbf{Fallback}: markiert den Schritt als \code{FAILED}.
        \item Der aufrufende Service bleibt erreichbar.
        \item Nach Wartezeit $\rightarrow$ Half-Open $\rightarrow$ Retry möglich.
      \end{itemize}

      \vspace{-0.1cm}
      \begin{block}{Wirkung}
        \small Isolation of Failures -- der Fehler bleibt lokal (ADR-04).
      \end{block}
    \end{column}
  \end{columns}
\end{frame}

\begin{frame}[fragile]{TA-04 -- Circuit Breaker: Konfiguration}
  \tabadge{TA-04}
  \begin{columns}[T,onlytextwidth]
    \begin{column}{0.5\textwidth}
      \begin{lstlisting}[style=yaml]
resilience4j:
  circuitbreaker:
    instances:
      nextService:
        sliding-window-size: 2
        minimum-number-of-calls: 1
        permitted-number-of-calls-in-half-open-state: 1
        failure-rate-threshold: 50
        wait-duration-in-open-state: 10s
        automatic-transition-from-open-to-half-open-enabled: true
      \end{lstlisting}
      \srcfile{services/fluid-analysis-service/src/main/resources/application.yml}
    \end{column}
    \begin{column}{0.46\textwidth}
      \textbf{Zustandsmodell}
      \begin{center}
      \begin{tikzpicture}[
        node distance=9mm,
        every node/.style={font=\scriptsize},
        st/.style={draw=darkblue,rounded corners,fill=lightblue,minimum width=15mm,minimum height=6mm,align=center},
        op/.style={draw=warnorange,rounded corners,fill=warnorange!10,minimum width=15mm,minimum height=6mm,align=center}
      ]
      \node[st] (c) {Closed};
      \node[op,right=of c] (o) {Open};
      \node[st,right=of o] (h) {Half-Open};
      \draw[-{Latex},thick] (c) -- node[above]{Fehler} (o);
      \draw[-{Latex},thick] (o) to[bend left=20] node[above]{10s} (h);
      \draw[-{Latex},thick] (h) to[bend left=20] node[below]{Fehler} (o);
      \draw[-{Latex},thick] (h) to[bend left=45] node[below]{Erfolg} (c);
      \end{tikzpicture}
      \end{center}

      \begin{block}{Health-Endpoint}
        \small \code{/actuator/health} zeigt den Breaker-Zustand.
      \end{block}
    \end{column}
  \end{columns}
\end{frame}

% --- TA-05 -----------------------------------------------------------------
\begin{frame}{TA-05 -- Apache Kafka (bewusst nicht umgesetzt)}
  \tabadge{TA-05}
  \begin{columns}[T,onlytextwidth]
    \begin{column}{0.52\textwidth}
      \textbf{Anforderung}
      \begin{itemize}
        \item Priorität \textbf{COULD} -- nicht Teil des Pflichtumfangs.
        \item Kafka könnte Statusnachrichten asynchron übertragen.
      \end{itemize}

      \vspace{0.3cm}
      \textbf{Entscheidung}
      \begin{itemize}
        \item Kein Kafka im Code.
        \item Statusmeldungen laufen synchron per REST-Callback (FR-09).
      \end{itemize}
    \end{column}
    \begin{column}{0.44\textwidth}
      \begin{block}{Begründung (ADR-02)}
        \small Für den kleinen PoC ist REST einfach nachvollziehbar, testbar und mit der Aufgabenstellung kompatibel. Kafka bleibt bewusst außerhalb des Kernumfangs.
      \end{block}
      \vspace{0.2cm}
      \begin{block}{Trade-off}
        \small REST erzeugt temporale Kopplung. Für eine produktive Variante wäre asynchrones Messaging zu prüfen.
      \end{block}
    \end{column}
  \end{columns}
\end{frame}

% ===========================================================================
% 04 - Architecture Decision Records (ADR) im Code
% ===========================================================================

% --- ADR-01 -----------------------------------------------------------------
\begin{frame}[fragile]{ADR-01 -- Fachlicher Schnitt statt technischer Layer}
  \adrbadge{ADR-01}
  \begin{columns}[T,onlytextwidth]
    \begin{column}{0.5\textwidth}
      \textbf{Entscheidung}
      \begin{itemize}
        \item Services werden nach \textbf{Business Capabilities} geschnitten.
        \item Kein technischer Layer-Schnitt (kein ``Controller-Service'').
      \end{itemize}

      \vspace{0.25cm}
      \textbf{Nachweis: Modulstruktur}
      \begin{lstlisting}[style=bash]
services/
  configuration-service/          -> Engine-Konfiguration
  analysis-management-service/    -> AnalysisRun, Status, Resultate
  fluid-analysis-service/         -> Oil-/Fuel-Analyse
  thermal-analysis-service/       -> thermische Analyse
  electrical-analysis-service/    -> elektrische Analyse
  engine-management-analysis-service/ -> EMS-Analyse
      \end{lstlisting}
    \end{column}
    \begin{column}{0.46\textwidth}
      \textbf{Interne Schichtung (je Service)}
      \begin{lstlisting}[style=bash]
fluid/
  api/             AnalysisController, AnalysisCommand
  application/     AnalysisWorker
  infrastructure/  NextServiceClient,
                   AnalysisManagementClient
  FluidAnalysisServiceApplication.java
      \end{lstlisting}

      \vspace{0.2cm}
      \begin{block}{Wirkung}
        \small Vermeidung des Anti-Patterns \textbf{Wrong Cut} (QR-06).
      \end{block}
    \end{column}
  \end{columns}
\end{frame}

% --- ADR-02 -----------------------------------------------------------------
\begin{frame}[fragile]{ADR-02 -- REST für die Kommunikation}
  \adrbadge{ADR-02}
  \begin{columns}[T,onlytextwidth]
    \begin{column}{0.52\textwidth}
      \begin{lstlisting}[style=java]
@RestController
@RequestMapping("/internal/analyses")
public class AnalysisController {

    private final AnalysisWorker worker;

    @PostMapping
    public ResponseEntity<Void> start(
            @RequestBody AnalysisCommand command) {
        worker.execute(command);
        return ResponseEntity.accepted().build();
    }
}
      \end{lstlisting}
      \srcfile{services/fluid-analysis-service/.../api/AnalysisController.java}
    \end{column}
    \begin{column}{0.44\textwidth}
      \textbf{Erläuterung}
      \begin{itemize}
        \item Alle Services kommunizieren über HTTP/REST mit definierten DTOs.
        \item Kein Service greift auf interne Klassen eines anderen zu.
        \item Ausgehende Aufrufe über \code{RestClient}.
      \end{itemize}
      \vspace{0.2cm}
      \begin{block}{Nachweis}
        \small \code{*Controller} (eingehend) + \code{RestClient}-Clients (ausgehend) in jedem Service.
      \end{block}
    \end{column}
  \end{columns}
\end{frame}

% --- ADR-03 -----------------------------------------------------------------
\begin{frame}[fragile]{ADR-03 -- Choreographie statt Orchestration}
  \adrbadge{ADR-03}
  \begin{columns}[T,onlytextwidth]
    \begin{column}{0.54\textwidth}
      \textbf{Start nur des Anchor-Algorithmus}
      \begin{lstlisting}[style=java]
public AnalysisRun start(String configurationId) {
    ConfigurationSnapshot configuration =
        configurationClient.get(configurationId);
    AnalysisRun run = AnalysisRun.start(
        "A-" + UUID.randomUUID(), configurationId);
    AlgorithmExecution fluid = run.execution(AlgorithmName.FLUID);
    fluid.updateStatus(AnalysisStatus.RUNNING, null);

    // Commit before external HTTP call (avoids callback race).
    run = repository.saveAndFlush(run);

    serviceStarter.start(AlgorithmName.FLUID,
        new AnalysisCommand(run.getAnalysisId(),
                            configuration, List.of()));
    return run;
}
      \end{lstlisting}
      \srcfile{.../application/AnalysisApplicationService.java}
    \end{column}
    \begin{column}{0.42\textwidth}
      \textbf{Weitergabe durch den Service}
      \begin{lstlisting}[style=java]
@Async
public void execute(AnalysisCommand command) {
    managementClient.reportStatus(
        command.analysisId(), ALGORITHM, "RUNNING", null);
    if (!pause(command.analysisId())) return;

    boolean ok = valid(command.configuration()
                       .get("oilSystem"))
              && valid(command.configuration()
                       .get("fuelSystem"));
    String result = ok ? "OK" : "FAILED";
    managementClient.reportResult(command.analysisId(),
        ALGORITHM, ok ? "READY" : "FAILED", result, null);

    if (ok) {
        nextServiceClient.startNext(command, result);
    }
}
      \end{lstlisting}
      \srcfile{.../application/AnalysisWorker.java}
    \end{column}
  \end{columns}
\end{frame}

% --- ADR-04 -----------------------------------------------------------------
\begin{frame}[fragile,shrink=5]{ADR-04 -- Circuit Breaker mit Resilience4j}
  \adrbadge{ADR-04}
  \begin{columns}[T,onlytextwidth]
    \begin{column}{0.54\textwidth}
      \begin{lstlisting}[style=java]
@CircuitBreaker(name = "analysisServiceStarter")
public void start(AlgorithmName algorithm,
                  AnalysisCommand command) {
    String baseUrl = switch (algorithm) {
        case FLUID -> fluidUrl;
        case THERMAL -> thermalUrl;
        case ELECTRICAL -> electricalUrl;
        case ENGINE_MANAGEMENT -> engineManagementUrl;
    };

    builder.baseUrl(baseUrl).build()
            .post().uri("/internal/analyses")
            .body(command).retrieve().toBodilessEntity();
}
      \end{lstlisting}
      \srcfile{.../infrastructure/AnalysisServiceStarter.java}
    \end{column}
    \begin{column}{0.42\textwidth}
      \textbf{Erläuterung}
      \begin{itemize}
        \item Umsetzung von \textbf{Isolation of Failures}.
        \item Der Fehler bleibt lokal im aufrufenden Service behandelbar.
        \item Zwei Breaker-Instanzen:
          \begin{itemize}
            \item \code{analysisServiceStarter} (Management $\rightarrow$ Analyse)
            \item \code{nextService} (Analyse $\rightarrow$ nächster Service)
          \end{itemize}
      \end{itemize}
      \begin{block}{Detail}
        \small Vollständiger Code + Fallback: siehe TA-04.
      \end{block}
    \end{column}
  \end{columns}
\end{frame}

% --- ADR-05 -----------------------------------------------------------------
\begin{frame}{ADR-05 -- Container pro Microservice}
  \adrbadge{ADR-05}
  \begin{columns}[T,onlytextwidth]
    \begin{column}{0.5\textwidth}
      \textbf{Entscheidung}
      \begin{itemize}
        \item Jeder Service besitzt ein eigenes \code{Dockerfile}.
        \item Jeder Service läuft in einem eigenen Container.
      \end{itemize}

      \vspace{0.2cm}
      \textbf{Nachweis}
      \begin{itemize}
        \item 6 $\times$ \code{Dockerfile} (Multi-Stage).
        \item \code{docker-compose.yml} mit 6 Services und Ports 8081--8086.
      \end{itemize}
    \end{column}
    \begin{column}{0.46\textwidth}
      \begin{center}
      \begin{tikzpicture}[
        every node/.style={font=\scriptsize},
        box/.style={draw=darkblue,rounded corners,fill=lightblue,minimum width=42mm,minimum height=6mm,align=center}
      ]
      \node[box] (a) {configuration-service :8081};
      \node[box,below=1.5mm of a] (b) {analysis-management :8082};
      \node[box,below=1.5mm of b] (c) {fluid-analysis :8083};
      \node[box,below=1.5mm of c] (d) {thermal-analysis :8084};
      \node[box,below=1.5mm of d] (e) {electrical-analysis :8085};
      \node[box,below=1.5mm of e] (f) {engine-management :8086};
      \end{tikzpicture}
      \end{center}

      \vspace{0.1cm}
      \begin{block}{Wirkung}
        \small Unterstützt Independent Deployability (QR-03).
      \end{block}
    \end{column}
  \end{columns}
\end{frame}

% --- ADR-06 -----------------------------------------------------------------
\begin{frame}[fragile,shrink=3]{ADR-06 -- Keine Shared Persistence}
  \adrbadge{ADR-06}
  \begin{columns}[T,onlytextwidth]
    \begin{column}{0.5\textwidth}
      \textbf{Getrennte Entities}
      \begin{lstlisting}[style=java]
@Entity
@Table(name = "engine_configurations")
public class EngineConfiguration {
    @Id
    private String configurationId;
    private String oilSystem;
    private String fuelSystem;
    // ...
}
      \end{lstlisting}
      \srcfile{configuration-service/.../domain/EngineConfiguration.java}

      \begin{lstlisting}[style=java]
@Entity
@Table(name = "analysis_runs")
public class AnalysisRun {
    @Id
    private String analysisId;
    // ...
}
      \end{lstlisting}
      \srcfile{analysis-management-service/.../domain/AnalysisRun.java}
    \end{column}
    \begin{column}{0.46\textwidth}
      \textbf{Integration über DTO}
      \begin{lstlisting}[style=java]
public record ConfigurationSnapshot(
        String configurationId,
        String oilSystem,
        String fuelSystem,
        String coolingSystem,
        String electricalSystem,
        String engineManagementSystem) {
}
      \end{lstlisting}
      \srcfile{.../infrastructure/ConfigurationSnapshot.java}
      \begin{block}{Nachweis stateless}
        \small Analyse-Services haben \textbf{keinen} \code{spring.datasource}-Block in \code{application.yml}.
      \end{block}
    \end{column}
  \end{columns}
\end{frame}

% ===========================================================================
% 05 - Live-Demo, FR/QR-Nachweise und Fazit
% ===========================================================================

\begin{frame}[fragile]{Live-Demo -- Circuit Breaker und Retry}
  \frbadge{FR-10/11}
  \begin{columns}[T,onlytextwidth]
    \begin{column}{0.5\textwidth}
      \textbf{Ablauf der Demo}
      \begin{enumerate}
        \item System starten:
          \begin{lstlisting}[style=bash]
docker compose up --build
          \end{lstlisting}
        \item Thermal stoppen:
          \begin{lstlisting}[style=bash]
docker compose stop thermal-analysis-service
          \end{lstlisting}
        \item Analyse starten $\rightarrow$ \code{THERMAL = FAILED}.
        \item Thermal starten:
          \begin{lstlisting}[style=bash]
docker compose start thermal-analysis-service
          \end{lstlisting}
        \item Retry ausführen $\rightarrow$ \code{overallResult = OK}.
      \end{enumerate}
    \end{column}
    \begin{column}{0.46\textwidth}
      \textbf{Retry-Endpoint}
      \begin{lstlisting}[style=bash]
POST /api/analyses/{analysisId}/algorithms/THERMAL/retry
      \end{lstlisting}

      \vspace{0.1cm}
      \textbf{Was die Demo zeigt}
      \begin{itemize}
        \item \textbf{TA-04 / ADR-04}: Circuit Breaker fängt den Ausfall ab.
        \item \textbf{FR-11}: System bleibt erreichbar, Schritt wird \code{FAILED}.
        \item \textbf{FR-10}: Nur der fehlgeschlagene Schritt wird wiederholt.
        \item \textbf{ADR-03}: Choreographie läuft ab Thermal weiter.
      \end{itemize}
    \end{column}
  \end{columns}
\end{frame}

\begin{frame}[fragile]{FR-07 -- Gesamtergebnis im Code}
  \frbadge{FR-07}
  \begin{columns}[T,onlytextwidth]
    \begin{column}{0.54\textwidth}
      \begin{lstlisting}[style=java]
public void recalculateOverallResult() {
    boolean failed = executions.stream()
        .anyMatch(e -> e.getStatus() == AnalysisStatus.FAILED
                    || e.getResult() == AnalysisResult.FAILED);

    if (failed) {
        overallResult = AnalysisResult.FAILED;
        return;
    }

    boolean unfinished = executions.stream()
        .anyMatch(e -> e.getStatus() == AnalysisStatus.PENDING
                    || e.getStatus() == AnalysisStatus.RUNNING);

    overallResult = unfinished ? null : AnalysisResult.OK;
}
      \end{lstlisting}
      \srcfile{analysis-management-service/.../domain/AnalysisRun.java}
    \end{column}
    \begin{column}{0.42\textwidth}
      \textbf{Regel (architecture.md \S8)}
      \begin{itemize}
        \item Ein Fehler $\rightarrow$ \code{FAILED}.
        \item Schritte offen $\rightarrow$ \code{null}.
        \item Alle vier fertig $\rightarrow$ \code{OK}.
      \end{itemize}
      \vspace{0.2cm}
      \begin{block}{Test}
        \small \code{AnalysisRunTest} prüft alle drei Fälle (QR-07).
      \end{block}
    \end{column}
  \end{columns}
\end{frame}

\begin{frame}[fragile]{QR-08 -- Responsiveness durch \code{@Async}}
  \qrbadge{QR-08}
  \begin{columns}[T,onlytextwidth]
    \begin{column}{0.5\textwidth}
      \begin{lstlisting}[style=java]
@PostMapping
public ResponseEntity<Void> start(
        @RequestBody AnalysisCommand command) {
    worker.execute(command);
    return ResponseEntity.accepted().build();
}
      \end{lstlisting}
      \srcfile{.../api/AnalysisController.java}

      \vspace{0.1cm}
      \begin{lstlisting}[style=java]
@Async
public void execute(AnalysisCommand command) {
    // laeuft im Hintergrund-Thread
}
      \end{lstlisting}
      \srcfile{.../application/AnalysisWorker.java}
    \end{column}
    \begin{column}{0.46\textwidth}
      \textbf{Erläuterung}
      \begin{itemize}
        \item Der Controller antwortet sofort mit \code{202 Accepted}.
        \item \code{worker.execute} läuft dank \code{@Async} im Hintergrund.
        \item Der Client blockiert nicht bis zum Ende der Analyse.
        \item Status ist währenddessen weiter abfragbar (QR-02).
      \end{itemize}
      \begin{block}{Zusammenspiel}
        \small \code{@EnableAsync} (TA-01) + \code{@Async} + \code{202 Accepted}.
      \end{block}
    \end{column}
  \end{columns}
\end{frame}

\begin{frame}{Fazit -- Anforderungen nachweisbar umgesetzt}
  \begin{columns}[T,onlytextwidth]
    \begin{column}{0.5\textwidth}
      \textbf{Technische Anforderungen}
      \begin{itemize}
        \item[\textcolor{okgreen}{\checkmark}] TA-01 Spring Boot
        \item[\textcolor{okgreen}{\checkmark}] TA-02 Docker
        \item[\textcolor{okgreen}{\checkmark}] TA-03 Docker Compose
        \item[\textcolor{okgreen}{\checkmark}] TA-04 Resilience4j
        \item[\textcolor{codegray}{--}] TA-05 Kafka (COULD, bewusst offen)
      \end{itemize}

      \vspace{0.3cm}
      \textbf{Architekturentscheidungen}
      \begin{itemize}
        \item[\textcolor{okgreen}{\checkmark}] ADR-01 bis ADR-06 im Code belegt
      \end{itemize}
    \end{column}
    \begin{column}{0.46\textwidth}
      \begin{block}{Kernbotschaft}
        \small Jede zentrale Anforderung ist über eine konkrete Datei und ein
        Symbol im Code nachweisbar -- nicht nur dokumentiert, sondern
        implementiert und testbar.
      \end{block}

      \vspace{0.2cm}
      \begin{block}{Referenzen}
        \small
        \code{docs/code-traceability.md}\\
        \code{docs/arc42.md}\\
        \code{docs/requirements.md}\\
        \code{scripts/e2e.sh}
      \end{block}
    \end{column}
  \end{columns}
\end{frame}

\begin{frame}[plain]
  \begin{center}
    \hbrslogo
    \vspace{1.5cm}

    {\Huge\bfseries\color{darkblue} Vielen Dank!}\par
    \vspace{0.6cm}
    {\large Fragen und Diskussion}\par

    \vspace{1.5cm}
    {\large Theodoros Bampis \& Yunus Emre Sirin}\par
    {\large SEKA -- Sommersemester 2026}\par
  \end{center}
\end{frame}


\end{document}
